\begin{abstract}
We develop a deterministic \emph{substrate theory fragment} in which observable physics is modeled as excitation dynamics of a three-dimensional brane-like medium embedded in $\R^4$, with external time as the evolution parameter.
The continuum description is written as a coordinate-free hyperelastic theory for an embedding field $\vecX: \Omega \times \R \to \R^4$, while a cubic lattice is treated as the candidate microphysical realization whose long-wavelength limit is the continuum model.
This separation is central: the continuum theory expresses the constitutive and geometric content cleanly, whereas the lattice supplies the actual branch structure, dispersion, and mixing mechanisms to be tested numerically.

Three mechanisms are isolated and formulated explicitly.
First, linearization about a homogeneous background yields polarized wave branches and an effective wave cone, making Lorentz-type kinematics interpretable as an emergent long-wavelength symmetry rather than a microscopic postulate.
Second, because transverse deformation changes intrinsic distances through the induced metric, homogeneous prestretch produces a transverse-to-lateral backreaction; in a weak and slowly varying regime this motivates a candidate contraction channel for gravity-like behavior.
Third, an ordered narrowband carrier sector supports adiabatic transport of polarization states and thereby induces Berry or Wilczek--Zee connections, which provide the natural geometric variables of a coarse-grained gauge description.

We also formulate localized particle-like states as nonlinear bound modes of the same substrate.
Here the emphasis is deliberately modest: the paper does not claim a derivation of the full Standard Model, QCD, or general relativity.
Its contribution is instead to connect the lattice postulate, continuum equations, geometric-phase diagnostics, and localized-mode program into one coherent mathematical framework with a clear computational agenda.
Within that agenda, the most informative near-term target is not an electron-like mode but a proton- or neutron-like \emph{baryon-scale triplet state}, because the cubic lattice naturally provides a three-channel axis-aligned sector whose mixing can be analyzed directly.
\end{abstract}
